% ----------------------------------------------------------
% Conclusão
% ----------------------------------------------------------
\chapter{Conclusão}

Este trabalho demonstrou a viabilidade de uma plataforma robótica seguidora de linha autônoma, modular e de código aberto, adaptada do iRobot Roomba® 614. A consolidação do Raspberry Pi 4 como único processador — executando ROS 2 Humble com \textit{Edge Computing} — eliminou microcontroladores intermediários, garantiu comunicação serial estável a 115200 bps e permitiu malha PID a 20 Hz com latência desprezível. A sintonia empírica ($K_p=0{,}8$, $K_i=0{,}0$, $K_d=0{,}2$) possibilitou seguimento preciso em curvas fechadas e retas, com desvios laterais zerados. A validação prévia em \textit{Turtlesim} confirmou a portabilidade da lógica de controle do simulador para o hardware físico. A plataforma serve como laboratório móvel para o ensino de controle automático, sistemas embarcados e arquiteturas distribuídas.

Como trabalhos futuros, destacam-se: (1)~integração física de câmera para ativar o módulo ArUco/OpenCV já validado em \textit{RViz2}, habilitando navegação por marcos visuais; (2)~implantação em cenários reais de logística \textit{indoor} (almoxarifados, unidades hospitalares); e (3)~coordenação de múltiplos robôs via DDS do ROS 2 para estudos de fluxo e evasão colaborativa de obstáculos.